% method: cong-thuc-poisson-nhiet
\begin{dieukienbox}
Dùng khi: phương trình \textbf{nhiệt} $u_t=k\Delta u$ trên toàn $\mathbb{R}^n$, dữ liệu $u(\cdot,0)=g$.
\end{dieukienbox}

\begin{nhobox}
\textbf{Công thức Poisson} (tích chập với nhân nhiệt $\Phi$):
\[
  u(x,t)=\int_{\mathbb{R}^n}\Phi(x-y,t)\,g(y)\,dy,\qquad
  \Phi(x,t)=\frac{1}{(4\pi k t)^{n/2}}\,e^{-|x|^2/(4kt)}.
\]
1D: $\displaystyle u(x,t)=\frac{1}{\sqrt{4\pi k t}}\int_{-\infty}^{\infty} e^{-(x-y)^2/(4kt)}\,g(y)\,dy.$
\end{nhobox}

\begin{meobox}
$g$ là hàm đặc trưng $\chi_{[a,b]}$ $\Rightarrow$ đổi biến $w=\frac{y-x}{\sqrt{4kt}}$, tích phân ra
hàm $\operatorname{erf}$ (xem Phụ lục B). Nhiều chiều: nhân nhiệt tách tích theo từng biến.
\end{meobox}